\section{Introduction}\label{sec:intro}

\subsection{The problem}\label{sec:intro-problem}

In the two-sided assignment game of buyers and sellers the core is never empty: the linear
programme behind the optimal assignment has an integral optimal solution, and its dual
solutions are the core payoffs \cite{ShapleyShubik71}. With three sides
this fails. A three-sided assignment game can
have an empty core, and the core is nonempty exactly when the fractional relaxation of the
three-dimensional assignment problem has an integral optimal solution. How often this happens
in a large random market has been an open question; the only figures found in the sources
we searched concern very
small markets, or markets with one small side, with independent team values
\cite{Lucas95}.

We study the following market. There are three sectors of $q$ agents. Every pair of agents
from different sectors has a value, and the $3q^2$ pair values are independent and uniform on
$[0,1]$. A team consists of one agent from each sector and is worth the sum of its three pair
values. A coalition is worth the best total value of disjoint teams that it can form. This is
the model of Frieze, Pegden and Tkocz \cite{FPT20}, read as a market: they showed that the
best assignment falls short of the trivial upper bound $3q$ by an amount of order $q^{1/3}$.

Let $\Tq$ be the value of the best perfect assignment, $\LPq$ the optimum of the fractional
relaxation, and $\Dq=\LPq-\Tq\ge0$ the gap. The core is nonempty if and only if $\Dq=0$
(Proposition~\ref{prop:model-core}). Full definitions are in Section~\ref{sec:model}.

\subsection{The main theorem}\label{sec:intro-main}

Throughout, $\betaint=2.302$ and $\alphalp=2.2519$.

\begin{theorem}\label{thm:main}
In the random three-sided market with independent uniform pair values, as $q\to\infty$:
\begin{enumerate}
\item[(i)] $\E\Dq\ \ge\ (\betaint-\alphalp-o(1))\,q^{1/3}\ \ge\ (0.050-o(1))\,q^{1/3}$; in
particular $\E\Dq\to\infty$.
\item[(ii)] $\Dq/\E\Dq\to1$ in probability.
\item[(iii)] $\P(\text{the core is nonempty})=\P(\Dq=0)\to0$.
\end{enumerate}
\end{theorem}

So the core of the random three-sided assignment market is empty with probability tending to
$1$. The theorem is asymptotic. The proof gives no bound at any size that can be simulated,
and the thresholds implicit in it are astronomically large (Section~\ref{sec:discussion}).
Two steps of the proof are computer-assisted (Sections~\ref{sec:intro-computer}
and~\ref{sec:certificates}).

\subsection{The four ingredients}\label{sec:intro-ingredients}

Write $\Phiq=3q-\Tq$ and $\PhiLP=3q-\LPq$ for the integer and the fractional deficit, so that
$\Dq=\Phiq-\PhiLP$. Both are of order $q^{1/3}$. The theorem follows from four statements.

\emph{The gap is asymptotically deterministic} (Theorem~\ref{thm:variance},
Section~\ref{sec:variance}). For every $\eps\in(0,\tfrac12)$, with high probability the
optimal assignment
uses no pair of deficit above $C_\eps q^{-1/2+\eps}$, and no team in the support of an optimal
fractional solution has deficit above $C_\eps q^{-1/2+\eps}$. With a Poincar\'e inequality this
gives $\Var\Dq=O(q^{-1/6+\eps})$. Hence $\P(\Dq=0)\to0$ as soon as $\E\Dq$ stays bounded away
from zero; far less than (i) is needed for (iii).

\emph{An integer lower bound} (Theorem~\ref{thm:integer}, Sections~\ref{sec:iid}
and~\ref{sec:transfer}). With high probability every perfect assignment has total deficit at
least $\betaint q^{1/3}$, that is, $\Phiq\ge\betaint q^{1/3}$. The bound is first proved in the
model with independent team deficits (Theorem~\ref{thm:iid}) and then carried to the market by
a comparison inequality (Theorem~\ref{thm:transfer}).

\emph{A fractional upper bound from the tree} (Theorem~\ref{thm:lp},
Sections~\ref{sec:localconv} and~\ref{sec:lp}). Let $\tree$ be the Poisson hypertree, the
local limit of the hypergraph of near-tight teams, and let $m$ be any invariant random perfect
matching on $\tree$, with $c(m)$ the expected cost of the root's team. Then
$\limsup_q\E\PhiLP/q^{1/3}\le c(m)$.

\emph{A cheap matching on the tree} (Theorems~\ref{thm:process} and~\ref{thm:existence},
Sections~\ref{sec:tree} and~\ref{sec:existence}). Let $(\xi_k)$ be a Poisson process of
intensity $\tfrac12\xi^2\dd\xi$ on $[0,\infty)$. A law $\mu$ on $\R$ is a fixed point of the
map $\RDE$ if
\begin{equation}\label{eq:intro-rde}
X\ \overset{d}{=}\ \min_k\bigl(\xi_k-X_k'-X_k''\bigr),
\end{equation}
where $X$ and all $X_k'$, $X_k''$ have law $\mu$ and are independent of each other and of
$(\xi_k)$. A fixed point defines an invariant perfect matching on $\tree$ of cost $9\,\E X$.
A fixed point exists, and every fixed point has $9\,\E X\le\alphalp$.

Given these, the proof of Theorem~\ref{thm:main} takes a few lines
(Section~\ref{sec:proofmain}). The last two ingredients give
$\E\PhiLP\le(\alphalp+o(1))q^{1/3}$. The second, with $\Phiq\ge0$, gives
$\E\Phiq\ge(\betaint-o(1))q^{1/3}$. Subtracting gives (i). Chebyshev's inequality with the
variance bound gives $\P(|\Dq-\E\Dq|>t\,\E\Dq)\le\Var\Dq/(t\,\E\Dq)^2\to0$ for every fixed
$t>0$, which is (ii), and (iii) is the case $t=1/2$.

\subsection{Overview of the proof}\label{sec:intro-overview}

The difficulty is that $\Phiq$ and $\PhiLP$ have the same order $q^{1/3}$ and, as far as the
numerical evidence goes, constants that differ by a few per cent. Three natural approaches
fail, each for a reason that shaped the proof.

\emph{The plain first moment is too weak.} There are $(q!)^2$ perfect assignments, and the
total deficit of a fixed one is a sum of $3q$ independent uniform variables. The expected
number of assignments with deficit at most $\beta q^{1/3}$ is therefore at most
$(q!)^2(\beta q^{1/3})^{3q}/(3q)!=(\beta^3e/27+o(1))^q$, which tends to zero only for
$\beta<3e^{-1/3}=2.1496\ldots$ This is a valid lower bound for $\Phiq/q^{1/3}$, but it lies
below $2.2495$, the cost of the matching on the tree constructed below, and below the
fractional optimum as measured numerically in the independent-team model (about $2.20$,
floating point). So it cannot separate the two optima.

\emph{No bounded structure carries the gap.} One might look for a local obstruction: a short
odd cycle in the support of the fractional optimum, a block of agents with a positive gap, a
local repair that rounds the fractional solution. But the neighbourhood of a typical agent in
the hypergraph of near-tight teams is a tree, and on the limiting tree the two problems have
the same value: every invariant fractional perfect matching on $\tree$ is the marginal of an
invariant random integral perfect matching of the same cost (Theorem~\ref{thm:decomp}). The
difference between $\LPq$ and $\Tq$ is that the finite market cannot realise the cheap
integral matchings of the tree, and this is not visible in any neighbourhood of bounded size.
In simulations the fractional optimum is fractional at most agents, so rounding it would need
of order $q$ repairs (Section~\ref{sec:discussion}).

\emph{The recursion is unstable under shifts.} The fixed-point equation \eqref{eq:intro-rde}
is the three-sided analogue of the recursion behind the limit of the random two-sided
assignment problem \cite{Aldous01}. Here, shifting the input law by a constant $c$ shifts
the output by $-2c$. A perturbation in the shift direction therefore doubles at every step,
plain iteration is unstable, and no narrow order interval of laws is invariant under $\RDE$.
Existence of a fixed point is not routine, and a numerical solution is not a proof: the
difference $\betaint-\alphalp$ is about $0.05$, and the cost $9\,\E X$ moves by $9c$ under a
shift by $c$.

The proof gets round these obstacles with three ideas.

\emph{A first moment conditioned on local stars} (Section~\ref{sec:iid}). In the
independent-team model, fix a truncation level $K$ and a partition of $[0,K]$ into bins, and
condition on the \emph{star} of every agent: the number of its teams in each bin. Given the
stars, the hypergraph of near-tight teams is uniform among simple hypergraphs with those
stars, which is bounded through the configuration model, and the expected number of
perfect assignments with a prescribed number of teams in each bin can be written down
exactly. When the stars are typical, its exponential rate is a concave functional of the usage
profile, and a tangent bound reduces the question to the sign of one explicit number
$\theta\beta+\Psi_0$. That number is negative for $\beta=2.30224$; this is the first
computer-assisted step. Conditioning matters because the unconditioned count is taken over all
star configurations at once, while the conditioned count is needed only for typical ones. The
bound is carried to the market by a comparison inequality (Section~\ref{sec:transfer}): the
expected number of cheap assignments in the market, restricted to a good event of the
near-tight hypergraph, is at most $e^{o(q)}$ times the same expectation in the
independent-team model on a slightly larger event, plus a negligible term. The comparison is
made for the conditional law of the near-tight teams outside a candidate assignment, given
the pair values of that assignment, on the event that few near-tight teams share a pair. The probability that a family of teams
contains no near-tight team is bounded by the argument of Boppana and Spencer
\cite{BoppanaSpencer89} for Janson's inequality, which uses Harris's inequality
\cite{Harris60} once.

\emph{Fractional solutions as conditional expectations of a tree process, with patching}
(Sections~\ref{sec:tree}, \ref{sec:localconv}, \ref{sec:lp}). A fixed point of
\eqref{eq:intro-rde} gives a family of messages on $\tree$, and with them an invariant random
perfect matching $m$ whose cost per team is $9\,\E X$ (Theorem~\ref{thm:process}). The
matching is not a function of the tree and its costs; it uses further randomness. What can be
computed from a finite neighbourhood is the conditional probability, given the ball of radius
$r$ around an agent, that a given team at that agent belongs to $m$. These conditional
probabilities are a fractional solution with small slack: their cost is $c(m)$ by the tower
property, and by martingale convergence the slack tends to zero as $r\to\infty$. Convergence
of marked balls in total variation (Theorem~\ref{thm:localconv}) carries the rule to the
finite market. The remaining slack is patched with an independent part of the randomness,
obtained by writing each pair deficit as the minimum of two independent variables; this
splitting device is Aldous's, from his work on the two-sided problem
\cite{Aldous92,Aldous01}. The cost of the patch is controlled by the upper bound of
\cite{FPT20}, Theorem~2(b), for the integer problem on the set of unserved agents. Averaging
over the patches gives an exact fractional assignment of cost
$(c(m)+o(1))q^{1/3}$. Fractionality enters at one place only, the conditional expectation.
The same construction would give an \emph{integral} assignment if $m$ were a measurable
function of the marked tree, and by the integer lower bound it cannot be one when
$c(m)<\betaint$ (Section~\ref{sec:discussion}).

\emph{The dispersive order} (Section~\ref{sec:existence}). Write a law as the law of $Q(E)$,
with $E$ a standard exponential variable and $Q$ nondecreasing, and say that
$Q_0\disp Q_1$ if $Q_1-Q_0$ is nondecreasing. This order ignores shifts, which removes the
unstable direction. On laws with increasing failure rate, the map $\RDE$ is monotone in this
order (Theorem~\ref{thm:mono}); the proof uses the theorem of Pr\'ekopa \cite{Prekopa73} that
integrals of log-concave functions are log-concave, and every image of $\RDE$ has increasing
failure rate. The iterates $\RDE^n(\delta_0)$ increase in the dispersive order, which
ignores shifts; normalised at a fixed level they converge, and the limit is a fixed point of
$\RDE$ after a shift that the shift law determines (Theorem~\ref{thm:existence}). A functional that equals the cost at a fixed point is antitone, so every fixed
point costs at most $c_n=3\bigl(\E\,\RDE(Q_n)+2\,\E\,Q_n\bigr)$ for every $n$, where $Q_n$ is
the $n$th iterate and $\E\,Q$ denotes the mean of the law $Q$. Monotonicity also makes a
one-sided numerical evaluation legitimate: replacing an iterate by a less dispersed law can
only increase the bound. Ten rounded iterations in interval arithmetic give
$9\,\E X\le\alphalp=2.2519$; this is the second computer-assisted step. In floating point the
cost of the fixed point is $2.2495456$.

\subsection{Related work}\label{sec:intro-related}

The statements in this subsection are limited to sources that were opened in full text during
the project, and say so where a source was not opened. The literature search was done by AI
agents; its queries and gaps are recorded with the project files. It did not cover Google
Scholar, SSRN or dissertation databases, and a paper with an unusual title could have been
missed.

\emph{The model.} Frieze, Pegden and Tkocz \cite{FPT20} study exactly the cost array
$W_{ijk}=a_{ij}+b_{ik}+c_{jk}$ with independent uniform entries, in minimisation form. Their
Theorem~2 shows that the expected optimum is at least $c_1n^{1/3}$ and that a greedy algorithm
achieves expected cost at most $c_2n^{1/3}$. The paper contains nothing on the fractional
relaxation, on duality or on the core. We use its Theorem~2(b) in the form: independent
uniform pair costs, bound in expectation. Several agents read this from the arXiv version;
the final text should quote the published statement. \todo{quote Theorem 2(b) of the
published version of FPT20 verbatim}

\emph{Random cores.} Lucas \cite{Lucas95} reports results of Crowley and Graham for
three-sided games with independent uniform values on triples, not pair-additive values: the
probability of a nonempty core is $209/240$ for two agents per side and about $73\%$ for
three, and it tends to $1$ when one side stays fixed and the others grow. The balanced limit
is not treated there. The primary sources behind these figures were not opened. For two
sides, Hassidim and Romm \cite{HassidimRomm15} prove an approximate law of one price in
random assignment games with bounded values; the core is always nonempty there, and the
question is its size. A related paper of Kanoria, Saban and Sethuraman on the size of the
core in two-sided markets \cite{KSS18} was opened in part. The papers \cite{KanekoWooders82} and \cite{Quint91} on
cores of partitioning and multi-sided games were not opened (the first at abstract level
only); no statement here rests on them.

\emph{Multi-dimensional assignment with independent costs.} Martin, M\'ezard and Rivoire
\cite{MMR05} treat costs that are independent over hyperedges by the cavity method. They
report that the replica-symmetric equations are inconsistent beyond two dimensions and do not
discuss a linear-programming relaxation. Frieze and Sorkin \cite{FriezeSorkin15} bound the
optimum for independent exponential costs and remark that even its scaling is not known. The
search found no source that compares the fractional and the integer optimum for a random
three-dimensional assignment problem.

\emph{Fractional and integral matchings in random hypergraphs.} At the level of thresholds
the two notions are close in the random $r$-uniform hypergraph. Kahn \cite{Kahn19} shows that
a perfect matching appears once the number of edges exceeds $(1+\eps)(n/r)\log n$. A result
of Krivelevich, known to us only as restated in \cite{DevlinKahn17} (the original was not
opened), says that at the first time without isolated vertices the random hypergraph has a
perfect fractional matching with high probability. Devlin and Kahn \cite{DevlinKahn17} study
perfect fractional matchings in $k$-out hypergraphs, including the $r$-partite case. These
are first-order statements about existence; our result concerns the constant on the scale
$q^{1/3}$ in a weighted problem.

\emph{The two-dimensional analogue.} For minimum matching on the complete graph with
exponential costs, Hessler and W\"astlund \cite{HesslerWastlund08} show that a relaxed
problem and the integer problem have expected costs within $O(\log n/n)$ of each other, less
than the standard deviation of either. The search found no result on the probability that the
fractional matching optimum is integral in that setting.

\emph{Methods.} The use of a limiting tree, recursive distributional equations and invariant
(unimodular) random networks follows the objective method
\cite{Aldous92,Aldous01,AldousSteele04,AldousLyons07}. The dispersive order is treated in
\cite{ShakedShanthikumar07}, and the variance bound uses the Poincar\'e inequality for the
exponential measure \cite{BobkovLedoux97}; see also \cite{BLM13}. Balanced matrices
\cite{FHO74} appear only in remarks. None of the works cited in this paragraph, nor
\cite{BoppanaSpencer89}, \cite{Harris60} and \cite{Prekopa73}, was opened during the project:
the attributions are from the agents' recollection, and the statements used from them are
either proved in the text or are standard. They must be checked against the sources before
any submission.

\subsection{What is computer-assisted}\label{sec:intro-computer}

Two inequalities between explicit real numbers are established by computation.

The first is $\Psi_0\le-3.84751328322$ for an explicit functional $\Psi_0$ of finitely many
parameters, which gives the constant $2.30224586\ge\betaint$ in Theorem~\ref{thm:iid} and,
with an explicit error budget, in Theorem~\ref{thm:integer}. The second is the bound
$9\,\E X\le2.2519$ of Theorem~\ref{thm:existence}, obtained by ten rounded iterations of the
map $\RDE$.

Both computations use interval arithmetic with outward rounding, and the second also uses
floating-point sums with a proved error bound. Their trusted base is the four basic
operations of the interval library and, for the second, IEEE arithmetic in the summations.
During the work the library's interval exponential and logarithm were found not to be strict
enclosures; neither certificate uses them. Section~\ref{sec:certificates} states exactly what
each program certifies, how to reproduce it, and which independent recomputations were made.
Everything else in the paper is a written proof. All numerical values labelled as floating
point, in remarks and in Section~\ref{sec:discussion}, are evidence and are not used in any
proof.
